% Generated by docs/generate_user_guide.py. Do not edit directly.
% Source: models registry

\begin{longtable}{@{}>{\raggedright\arraybackslash}p{0.16\textwidth}>{\raggedright\arraybackslash}p{0.14\textwidth}>{\raggedright\arraybackslash}p{0.36\textwidth}>{\raggedright\arraybackslash}p{0.26\textwidth}@{}}
\toprule
\textbf{Model} & \textbf{Used in} & \textbf{Functional form} & \textbf{Notes} \\
\midrule
\endfirsthead
\toprule
\textbf{Model} & \textbf{Used in} & \textbf{Functional form} & \textbf{Notes} \\
\midrule
\endhead
Double Exponential & Approach 1, Approach 2 & \(y = a e^{b x} + c e^{d x}\) & Sum of two exponential terms. It can represent profiles with two apparent rates, but the added flexibility may require more data to estimate stably. \\
Exponential & Approach 1, Approach 2 & \(y = a e^{b x}\) & Two-parameter exponential relationship. It can describe monotonic change that accelerates or decelerates exponentially over the fitted range. \\
Linear & Approach 1 & \(y = a + b x\) & Straight-line empirical relationship with an intercept and slope. It is appropriate when the in vitro/in vivo relationship is approximately linear over the fitted range. \\
Linear (zero intercept) & Approach 1 & \(y = a x\) & Straight-line empirical relationship constrained to pass through zero. It may be useful when a zero in vitro value should correspond to a zero in vivo value under the selected preprocessing and response scale. \\
Mass Loss (Gompertz) & Approach 2 & \(y = y_0 - A \exp\{-\exp[(e r_{max}/A)(t_{lag}-x)+1]\}\) & Mass-remaining model based on the modified Gompertz curve. It is intended for asymmetric sigmoidal mass-loss profiles with an apparent lag or plateau, a maximum loss-rate region, and a terminal loss extent. \\
Mass Loss (One-Phase) & Approach 2 & \(y = y_0 - A/(1 + e^{-(x-t_{50})/s})\) & Mass-remaining model with a delayed sigmoidal loss phase. It is intended for mass-loss profiles that begin constant and show a later drop. \\
Mass Loss (Surface Erosion) & Approach 2 & \(y = y_0 \max(1-kx,0)^n\) & Mass-remaining model for surface-controlled erosion. The fitted exponent n allows the profile shape to adapt to slab-like, cylindrical, spherical, or empirical geometry-like erosion behavior. \\
Mass Loss (Two-Phase) & Approach 2 & \(y = y_0 - B(1-e^{-k_b x}) - D/(1+e^{-(x-t_{50})/s})\) & Mass-remaining model with an optional early burst-loss phase followed by a delayed sigmoidal loss phase. It is intended for mass-loss profiles that may show an initial drop, an intermediate plateau, and a later drop. \\
Power & Approach 1, Approach 2 & \(y = a x^{b}\) & Power-law empirical relationship. It can describe nonlinear monotonic relationships where the response changes as a fitted power of time or the mapped predictor. \\
Stretched Exponential & Approach 2 & \(y = a e^{b x^{c}}\) & Exponential relationship with a fitted stretching exponent. It can describe non-single-rate behavior, but the additional shape parameter can increase sensitivity when data are sparse. \\
\bottomrule
\end{longtable}
